% Generated draft: requires mathematical and source review.
\documentclass[11pt]{article}
\usepackage[margin=25mm]{geometry}
\usepackage{fontspec}
\setmainfont{Latin Modern Roman}
\usepackage{amsmath,amssymb,mathtools,unicode-math}
\setmathfont{Latin Modern Math}
\usepackage{enumitem,needspace,xurl,hyperref}
\setlength{\parindent}{0pt}
\setlength{\parskip}{5pt}
\setlength{\emergencystretch}{3em}
\newcommand{\N}{\mathbb{N}}
\newcommand{\Z}{\mathbb{Z}}
\newcommand{\Q}{\mathbb{Q}}
\newcommand{\R}{\mathbb{R}}
\newcommand{\C}{\mathbb{C}}
\newcommand{\F}{\mathbb{F}}
\newcommand{\A}{\mathbb{A}}
\newcommand{\PP}{\mathbb{P}}
\newcommand{\E}{\mathbb{E}}
\newcommand{\eps}{\varepsilon}
\newcommand{\dd}{\,\mathrm{d}}
\begin{document}
\section*{UnsolvedMath ID: 30006659}
\textbf{Source problem code:} catalogue entry
\subsection{Definitions and mathematical statement}
For a prime \(p\), an elementary abelian \(p\)-group of rank \(k\) is a group isomorphic to \((\Z/p\Z)^k\). An action of a finite group \(G\) on a topological space \(X\) is free if \(gx=x\) implies \(g=1\) for every \(g\in G\) and \(x\in X\). Actions are interpreted here as continuous actions. The notation \(X\simeq Y\) denotes homotopy equivalence.

The supplied reconstructed statement asks for the following characterization, attributed in the supplied curation to Benson--Carlson. Under the explicit provisional convention that the number of sphere factors and all sphere dimensions are positive integers, for every finite group \(G\) and every integer \(r\geq 1\), prove that
\[
\begin{gathered}
\exists\, n_1,\ldots,n_r\in\Z_{\geq 1}\;\exists\, X\text{ a finite CW complex}:\\
X\simeq\prod_{i=1}^{r}S^{n_i}
\quad\text{and}\quad G\text{ acts freely on }X
\end{gathered}
\quad\Longleftrightarrow\quad
\begin{gathered}
\text{for every prime }p,\text{ there is no subgroup}\\
E\leq G\text{ with }E\cong(\Z/p\Z)^{r+1}.
\end{gathered}
\]
The source does not specify the allowed sphere dimensions or the range of \(r\); consequently, these positivity conventions require confirmation. The dimensions in this display are existentially quantified and may depend on \(G\) and \(r\).

The particular question is whether, for every odd prime \(p\), there exist positive integers \(n,m\) and a finite CW complex \(X\simeq S^n\times S^m\) admitting a free action of \(Qd(p)\). The supplied record does not define \(Qd(p)\). If the conventional notation is intended, it means
\[
Qd(p)=(\F_p)^2\rtimes\operatorname{SL}_2(\F_p),
\]
where \(\F_p\) is the field with \(p\) elements, \((\F_p)^2\) is its additive vector group, and \(\operatorname{SL}_2(\F_p)\) acts by its natural matrix action. This identification, as well as the positivity convention for \(n,m\), requires source confirmation.

The supplied original excerpt asks which finite groups act freely on a \emph{given} product of spheres. This differs from the reconstructed statement, which permits an arbitrary finite CW complex of the indicated homotopy type and existentially chosen dimensions. The precise relationship between these targets is unresolved by the supplied excerpt.
\subsection{Short English statement}
The reconstructed question asks whether excluding elementary abelian subgroups of rank \(r+1\), at every prime, exactly characterizes finite groups that act freely on some finite CW complex homotopy equivalent to a product of \(r\) spheres. It singles out \(Qd(p)\) for odd primes and two sphere factors. Sphere dimensions, the meaning of \(Qd(p)\), and the difference from the original question about a given product require confirmation.
\subsection{Sources}
\begin{itemize}
\item Local UnsolvedMath registry, permanent ID \(30006659\), title \emph{Free Finite-Group Actions on Products of Spheres}, source problem code \texttt{catalogue entry}. The supplied project identifies \texttt{data/unsolvedmath/problems.json} and the supplemental \texttt{catalog/problems.json} as local registries; no file lookup was performed.
\item \emph{Cohomology of Finite Groups: Interactions and Applications} (2026), Oberwolfach Reports, DOI \texttt{10.4171/owr/2026/9}, pp.\ 541--610. Supplied report reference: \url{https://doi.org/10.4171/owr/2026/9}. The original excerpt names the contribution \emph{Free symmetries of products of spheres and Steenrod closed parameter ideals}, by Marc Stephan, joint work with Henrik R\"uping and Erg\"un Yal\c{c}{\i}n. The record does not identify a separate problem proposer.
\item Adem--Smith, \emph{Periodic Complexes and Group Actions}: \url{https://arxiv.org/abs/math/0010096}.
\item Jackson, \emph{Qd(p)-free rank two groups}: \url{https://arxiv.org/abs/math/0503746}.
\item Supplied literature assessment, checked \texttt{2026-08-22}: the record labels the problem partially solved and states that the rank-one case and substantial rank-two families are solved, while the Benson--Carlson characterization still reduces to unresolved free actions for \(Qd(p)\) at odd primes. This is supplied dated curation derived from the OWR record, not an independently verified literature assessment; the cited claims and time-sensitive status require independent checking.
\item The supplied statement status is \texttt{corrected\_verified}, with the assertion that the statement was reconstructed and checked against the official Oberwolfach report. This is a supplied verification claim, not verification performed for this draft.
\end{itemize}
\end{document}
